% ============================================================
% Proof-writing macros
% ============================================================
\newcommand{\Given}{\par\noindent\textbf{Given. }\ignorespaces}
\newcommand{\Goal}{\par\noindent\textbf{Goal. }\ignorespaces}
\newcommand{\Strategy}{\par\noindent\textbf{Strategy. }\ignorespaces}
\newcommand{\ProofPlan} {\textbf{Plan.}\ }
\newcommand{\Recall}    {\textbf{Recall.}\ }
\newcommand{\Claim}     {\textbf{Claim.}\ }   % inline bold marker; \begin{claim} = numbered environment
\newcommand{\Wlog}      {\textit{Without loss of generality}}
\newcommand{\Fix}       {\textbf{Fix.}\ }
\newcommand{\Let}       {\textbf{Let.}\ }
\newcommand{\Choose}    {\textbf{Choose.}\ }
\newcommand{\Take}      {\textbf{Take.}\ }
\newcommand{\Note}      {\textbf{Note.}\ }
\newcommand{\ByDef}     {\textbf{By definition,}\ }
\newcommand{\ByAss}     {\textbf{By assumption,}\ }
\newcommand{\ByThm}[1]  {\textbf{By #1,}\ }
\newcommand{\Thus}      {\textbf{Thus,}\ }
\newcommand{\Hence}     {\textbf{Hence,}\ }
\newcommand{\Therefore} {\textbf{Therefore,}\ }
\newcommand{\AsReq}     {\textbf{as required.}}
\newcommand{\checkbox}  {$\square$\ }
\newcommand{\pp}        {\texttt{++}}
\providecommand{\conv}{\operatorname{conv}}
\providecommand{\argmin}{\operatorname*{arg\,min}}
\providecommand{\argmax}{\operatorname*{arg\,max}}

% ============================================================
% Nested numbered claims (reset per proof)
%   Claim 1, Claim 2, ...
%   Subclaim 1.1, 1.2, ...
% ============================================================
\newcounter{claim}
\newcounter{subclaim}[claim]

\renewcommand{\theclaim}{\arabic{claim}}
\renewcommand{\thesubclaim}{\theclaim.\arabic{subclaim}}

\AtBeginEnvironment{proof}{%
  \setcounter{claim}{0}%
  \setcounter{subclaim}{0}%
}

\newenvironment{claim}{%
  \refstepcounter{claim}%
  \par\medskip
  \noindent\textbf{Claim \theclaim.}\ \itshape
}{%
  \par\medskip
}

\newenvironment{subclaim}{%
  \refstepcounter{subclaim}%
  \par\medskip
  \noindent\textbf{Subclaim \thesubclaim.}\ \itshape
}{%
  \par\medskip
}

\newcommand{\ProofStep}[2]{\medskip\noindent\textbf{Step #1\,—\,#2}\quad}
% ============================================================
% Proof annotation tags — FIXED VERSION
% ============================================================
%
% Problem: \marginpar is forbidden inside tabular, tcolorbox, theorem
% environments, and other restricted modes. Calling \prooftag in those
% contexts raises the @marbox / \@marbox undefined error.
%
% Solution: Two separate primitives.
%
%   \prooftag{color}{label}   — margin version (use in proof body prose)
%   \inlinetag{color}{label}  — inline colored badge (use inside tabular T column,
%                               tcolorbox, definition, theorem environments)
%
% All named tag macros (e.g. \tagDU) produce the MARGIN version by default.
% Inside a tabular or box, call the \itag* inline variants instead.
% ============================================================

% ---- Base primitives ----------------------------------------

% Proof tag: inline colored badge.
% Originally used \marginpar but replaced with inline to avoid
% "Floats lost" errors at \clearpage in long documents.
\newcommand{\prooftag}[2]{%
  \textbf{\textcolor{#1}{\scriptsize #2}}%
}

% Inline tag: colored bold badge, zero extra vertical space.
% Safe everywhere: tabular, tcolorbox, theorem bodies, definition boxes.
\newcommand{\inlinetag}[2]{%
  \textbf{\textcolor{#1}{\scriptsize #2}}%
}

% ---- Core structural tags -----------------------------------
%   \tagXX    = margin version  (default, proof prose)
%   \itagXX   = inline version  (tabular T column, boxes)

% DU = Definition Unpacked
\newcommand{\tagDU} {\prooftag{blue!70!black}{DU}}
\newcommand{\itagDU}{\inlinetag{blue!70!black}{DU}}

% TA = Theorem / Axiom Applied
\newcommand{\tagTA} {\prooftag{green!50!black}{TA}}
\newcommand{\itagTA}{\inlinetag{green!50!black}{TA}}

% CS = Case Split
\newcommand{\tagCS} {\prooftag{red!60!black}{CS}}
\newcommand{\itagCS}{\inlinetag{red!60!black}{CS}}

% EX = Existence / witness choice
\newcommand{\tagEX} {\prooftag{magenta!60!black}{EX}}
\newcommand{\itagEX}{\inlinetag{magenta!60!black}{EX}}

% NEG = Negation unpacked (contradiction / contrapositive)
\newcommand{\tagNEG} {\prooftag{black}{NEG}}
\newcommand{\itagNEG}{\inlinetag{black}{NEG}}

% ---- ε–N / tail-logic tags ----------------------------------

% AM = Algebraic Manipulation
\newcommand{\tagAM} {\prooftag{black}{AM}}
\newcommand{\itagAM}{\inlinetag{black}{AM}}

% EA = ε arbitrary (start of universal ε argument)
\newcommand{\tagEA} {\prooftag{black}{EA}}
\newcommand{\itagEA}{\inlinetag{black}{EA}}

% NE = choose N = N(ε)
\newcommand{\tagNE} {\prooftag{black}{NE}}
\newcommand{\itagNE}{\inlinetag{black}{NE}}

% TN = take tail indices ≥ N
\newcommand{\tagTN} {\prooftag{black}{TN}}
\newcommand{\itagTN}{\inlinetag{black}{TN}}

% TR = triangle inequality / norm inequality
\newcommand{\tagTR} {\prooftag{black}{TR}}
\newcommand{\itagTR}{\inlinetag{black}{TR}}

% ---- Completeness / extraction tags -------------------------

% LUB = least-upper-bound / sup step
\newcommand{\tagLUB} {\prooftag{black}{LUB}}
\newcommand{\itagLUB}{\inlinetag{black}{LUB}}

% GLB = greatest-lower-bound / inf step
\newcommand{\tagGLB} {\prooftag{black}{GLB}}
\newcommand{\itagGLB}{\inlinetag{black}{GLB}}

% SS = subsequence selection / extraction
\newcommand{\tagSS} {\prooftag{black}{SS}}
\newcommand{\itagSS}{\inlinetag{black}{SS}}

% ---- Induction tags -----------------------------------------

% IH = inductive hypothesis applied
\newcommand{\tagIH} {\prooftag{purple!70!black}{IH}}
\newcommand{\itagIH}{\inlinetag{purple!70!black}{IH}}

% BC = base case
\newcommand{\tagBC} {\prooftag{teal!70!black}{BC}}
\newcommand{\itagBC}{\inlinetag{teal!70!black}{BC}}

% IS = inductive step
\newcommand{\tagIS} {\prooftag{teal!50!black}{IS}}
\newcommand{\itagIS}{\inlinetag{teal!50!black}{IS}}

% ------------------------------------------------------------
% Removed on purpose (not in ideal minimal set)
% ------------------------------------------------------------
% AM (algebraic manipulation): too granular; prose should carry it.
% UW (universal/witness intro): redundant with TN/EX in practice.
% ------------------------------------------------------------

% ============================================================
% Relation symbols
% ============================================================
\newcommand{\colonequiv}{\vcentcolon\equiv}   % predicate definitional equivalence  P :≡ Q
\newcommand{\coloniff}{\;\vcentcolon\iff\;}   % biconditional form of the above

% ============================================================
% Tolerance / Proximity quantifier macros
%
% Usage in predicate readings:
%
%   \ForallOT{\varepsilon}   --  \forall OutputTolerance(\varepsilon)
%   \ExistsIT{\delta}        --  \exists InputTolerance(\delta)
%   \ForallIT{\delta}        --  \forall InputTolerance(\delta)
%                                (used in uniform continuity: both
%                                 x and u range over inputs)
%
% These make the quantifier-typed reading of epsilon-delta
% definitions compact and semantically explicit.
%
% Full forms (for display when space allows):
%   \forall\,\operatorname{OutputTolerance}(\varepsilon)
%   \exists\,\operatorname{InputTolerance}(\delta)
% ============================================================
\newcommand{\ForallOT}[1]{%
  \forall\,\operatorname{OutputTolerance}(#1)}
\newcommand{\ExistsIT}[1]{%
  \exists\,\operatorname{InputTolerance}(#1)}
\newcommand{\ForallIT}[1]{%
  \forall\,\operatorname{InputTolerance}(#1)}

% ============================================================
% Math operators
% ============================================================
\DeclareMathOperator{\id}{id}
\DeclareMathOperator{\dom}{dom}
\DeclareMathOperator{\cod}{cod}
\DeclareMathOperator{\im}{im}

% ============================================================
% Citation macros
% ============================================================
\newcommand{\srccite}[1]{%
  \textsc{\citeauthor{#1}}%
}
\newcommand{\srcciteauthor}[1]{%
  \textsc{\citeauthor{#1}}%
}
\newcommand{\srccitesection}[2]{%
  \textsc{\citeauthor{#1}}, \textit{#2} [\citeyear{#1}]%
}

% ============================================================
% Proof vault links
% ============================================================
\newcommand{\ProofVaultURL}[1]{%
  \begin{remark*}[Proof vault]
  \href{#1}{Handwritten proof record}
  \end{remark*}%
}

% ============================================================
% Structural metadata markers — no PDF output
% Extraction markers: signature, structure, classification, proof-vault-url
% ============================================================
\newcommand{\LRAProofFor}[1]{}            % proof-file target theorem/proposition/lemma/corollary label
\newcommand{\LRAProofBodyStart}{\mbox{}\par\nobreak\noindent\ignorespaces}
\newcommand{\SignatureMetadata}[1]{}       % future Knowledge Explorer signature node metadata
\newcommand{\StructureMetadata}[1]{}       % future Knowledge Explorer structure blueprint metadata
\newcommand{\ClassificationMetadata}[1]{}  % future Knowledge Explorer classification metadata
\newcommand{\SyntacticMetadata}[1]{}       % future Knowledge Explorer syntactic object metadata

% ============================================================
% Print edition routing
% ============================================================
% Default: complete digital/PDF edition.
% Print edition may be enabled either by:
%   LRA_PRINT_EDITION=1 latexmk -lualatex main.tex
% or by command-line TeX injection:
%   lualatex "\def\LRAPrintEditionRequested{1}% =========================================================
% main.tex - Standalone root for Volume VIII
%
% Build (local): latexmk -lualatex main.tex
% Build (Overleaf): Menu -> Main Document -> main.tex
% =========================================================
\documentclass[12pt, a4paper]{report}
% =========================================================
% common/volume-preamble.tex
% Shared preamble for standalone per-volume root files.
%
% Usage --- in any volume-N-main.tex:
%   \documentclass[12pt, a4paper]{report}
%   % =========================================================
% common/volume-preamble.tex
% Shared preamble for standalone per-volume root files.
%
% Usage --- in any volume-N-main.tex:
%   \documentclass[12pt, a4paper]{report}
%   \input{common/volume-preamble}
%   \begin{document}
%   ...
%   \end{document}
%
% Rule: keep this in sync with the preamble block in main.tex.
% If you add a package to main.tex, add it here too.
% =========================================================

% -- Shared infrastructure ---------------------------------
\input{common/preamble}
\input{common/boxes}
\input{common/breadcrumb-macros}
\input{common/colors}
\input{common/structural-presentations}
\input{common/environments}
\input{common/macros}

% -- Hyperlinks and cross-references -----------------------
\usepackage{hyperref}
\usepackage[nameinlink,capitalize,noabbrev]{cleveref}
\usepackage{bookmark}

% -- cleveref configuration --------------------------------
\crefname{theorem}{Theorem}{Theorems}
\Crefname{theorem}{Theorem}{Theorems}
\crefname{claim}{Claim}{Claims}
\Crefname{claim}{Claim}{Claims}
\crefname{subclaim}{Subclaim}{Subclaims}
\Crefname{subclaim}{Subclaim}{Subclaims}

\hypersetup{
  colorlinks = true,
  linkcolor  = blue!60!black,
  citecolor  = blue!60!black,
  urlcolor   = blue!60!black,
}

% Allow \hyperref inside theorem optional arguments.
\pdfstringdefDisableCommands{%
  \def\hyperref[#1]#2{#2}%
}

% -- Indexes -----------------------------------------------
\usepackage{imakeidx}
\makeindex[name=tech,title={Technique Index}]
\makeindex[name=dep,title={Dependency Index}]

% -- Numbering and TOC depth -------------------------------
\setcounter{secnumdepth}{3}
\setcounter{tocdepth}{3}

%   \begin{document}
%   ...
%   \end{document}
%
% Rule: keep this in sync with the preamble block in main.tex.
% If you add a package to main.tex, add it here too.
% =========================================================

% -- Shared infrastructure ---------------------------------
% ============================================================
% Page layout
% ============================================================
\usepackage[margin=1in]{geometry}

% Make text that runs past the page margin visible in the PDF.
% LaTeX already reports these as "Overfull \hbox" warnings in the build log.
\hfuzz=0pt
\overfullrule=5pt

% ============================================================
% Fonts and encoding
% ============================================================
\usepackage[T1]{fontenc}

% ============================================================
% AMS mathematics
% ============================================================
\usepackage{amsmath, amssymb, amsthm}
\usepackage{mathtools}
\usepackage{stmaryrd}

% ============================================================
% Lists and tables
% ============================================================
\usepackage{array}
\usepackage{booktabs}
\usepackage{enumitem}
\usepackage{longtable}
\usepackage{makecell}
\usepackage{tabularx}
\usepackage{cancel}

% ============================================================
% Graphics (images)
% ============================================================
\usepackage{graphicx}
\usepackage{caption}


\usepackage{centernot}

% ============================================================
% Colours, boxes, and links
% ============================================================
\usepackage{xcolor}
\usepackage{mdframed}
\usepackage{tcolorbox}
\tcbuselibrary{skins,breakable}

% ============================================================
% TikZ
% ============================================================
\usepackage{tikz}
\usepackage{tikz-3dplot}
\usetikzlibrary{
  arrows.meta,
  backgrounds,
  calc,
  decorations.pathreplacing,
  decorations.pathmorphing,
  decorations.markings,
  fit,
  patterns,
  positioning,
  shapes.geometric
}

% ============================================================
% Miscellaneous
% ============================================================
\usepackage{etoolbox}
\usepackage{expl3}
\usepackage{float}
\usepackage{microtype}
\usepackage{multicol}
\usepackage{placeins}
\usepackage{xparse}
\usepackage{multirow}
\usepackage{rotating}   % sidewaysfigure

% ============================================================
% Bibliography
% ============================================================
\usepackage[authoryear]{natbib}



\usepackage{imakeidx}
\makeindex[name=tech,title={Technique Index}]
\makeindex[name=dep,title={Dependency Index}]

\input{common/exercise-format}




% --- Blueprint macros (merged from blueprint-macros.tex) ---
% ============================================================
%  Blueprint shared macros — \input{blueprint-macros} once
%  after \begin{document} before any blueprint \input calls.
%
%  Required in main.tex preamble:
%    \usepackage{tcolorbox}
%    \tcbuselibrary{skins}
%    \usepackage{amsthm,amsmath,amssymb,xcolor}
%    all \definecolor{cSet}...{cVS} definitions
%    \newtheorem{definition}{Definition}[chapter]  (or similar)
%    \newtheorem{remark}{Remark}[chapter]
% ============================================================

% ── Section boxes ─────────────────────────────────────────────
% NEW: strong frame in structure colour
\newtcolorbox{bpnew}[2]{
  enhanced,
  colback   = #1!14!white,
  colframe  = #1!80!black,
  coltitle  = white,
  fonttitle = \bfseries\scshape\small,
  title     = #2,
  boxrule   = 1.0pt, arc=2pt,
  left=7pt, right=7pt, top=4pt, bottom=5pt,
  toptitle=3pt, bottomtitle=3pt,
  before skip=0pt, after skip=7pt,
}

% INHERITED: faded frame in ancestor colour
\newtcolorbox{bpinh}[2]{
  enhanced,
  colback   = #1!7!white,
  colframe  = #1!38!white,
  coltitle  = #1!70!black,
  fonttitle = \bfseries\scshape\small,
  title     = #2,
  boxrule   = 0.6pt, arc=2pt,
  left=7pt, right=7pt, top=4pt, bottom=5pt,
  toptitle=3pt, bottomtitle=3pt,
  before skip=0pt, after skip=7pt,
}

% ── Entry line ────────────────────────────────────────────────
% \bpitem{symbol}{math}{gloss}   (gloss may be empty)
\newcommand{\bpitem}[3]{%
  \noindent
  {\small\bfseries #1}\hspace{6pt}%
  $\displaystyle #2$%
  \ifx\relax#3\relax\else
    \hfill{\small\rmfamily\itshape\color{lbltext!75}\,[#3]}%
  \fi
  \par\vspace{3pt}%
}

% ── Page header ───────────────────────────────────────────────
% \bpheader{colour}{Name}{signature}
\newcommand{\bpheader}[3]{%
  \noindent\vbox{\offinterlineskip
    \hbox{\colorbox{#1!85!black}{%
      \parbox[c][2.4em]{\dimexpr\linewidth-2\fboxsep\relax}{%
        \color{nodetext}%
        \hspace{8pt}{\large\bfseries #2}%
        \enspace{\normalsize\ttfamily\color{sigtext} #3}%
        \hfill{\footnotesize\scshape\bfseries Structure Blueprint\hspace{8pt}}%
      }%
    }}%
  }\par\vspace{8pt}%
}

% ── Footer ────────────────────────────────────────────────────
% \bpfooter{left}{right}
\newcommand{\bpfooter}[2]{%
  \vfill
  \noindent\hrule\vspace{4pt}%
  \noindent
  \begin{minipage}[t]{0.49\linewidth}\scriptsize #1\end{minipage}%
  \hfill
  \begin{minipage}[t]{0.49\linewidth}\scriptsize #2\end{minipage}%
}

% =========================================================
% Dependency Figure Styles
% Add to common/boxes.tex
% =========================================================

% Requires TikZ libraries:
% \usetikzlibrary{arrows.meta,positioning,calc}

\tikzset{
  depaxiom/.style={
    draw,
    rounded corners=3pt,
    fill=axiombox,
    draw=axiomborder,
    minimum width=4.2cm,
    minimum height=0.82cm,
    align=center,
    font=\small
  },
  depdefinition/.style={
    draw,
    rounded corners=3pt,
    fill=defbox,
    draw=defborder,
    minimum width=4.2cm,
    minimum height=0.82cm,
    align=center,
    font=\small
  },
  deptheorem/.style={
    draw,
    rounded corners=3pt,
    fill=thmbox,
    draw=thmborder,
    minimum width=4.2cm,
    minimum height=0.82cm,
    align=center,
    font=\small
  },
  depneutral/.style={
    draw,
    rounded corners=3pt,
    fill=gray!8,
    draw=gray!65,
    minimum width=4.2cm,
    minimum height=0.82cm,
    align=center,
    font=\small
  },
  depimplies/.style={->, thick, blue!70!black, >=Latex},
  depuses/.style={->, thick, yellow!60!black, >=Latex},
  depdefines/.style={->, thick, green!50!black, >=Latex},
depedgelab/.style={
  font=\tiny,
  fill=white,
  inner sep=0.5pt,
  text=black
},
  deplegendbox/.style={
    draw,
    rounded corners=3pt,
    fill=white,
    draw=gray!55,
    inner sep=5pt
  }
}

\newcommand{\DependencyLegend}{%
\begin{tikzpicture}[baseline=(current bounding box.north)]
\node[deplegendbox, anchor=north east] (L) {
\begin{tabular}{@{}ll@{}}
\multicolumn{2}{@{}l@{}}{\textbf{\small Legend}}\\[0.2em]
\tikz{\node[depaxiom, minimum width=0.7cm, minimum height=0.28cm] {};} & \scriptsize Axiom\\
\tikz{\node[depdefinition, minimum width=0.9cm, minimum height=0.35cm] {};} & \scriptsize Definition\\
\tikz{\node[deptheorem, minimum width=0.9cm, minimum height=0.35cm] {};} & \scriptsize Theorem / Proposition / Corollary\\
\tikz{\draw[depimplies] (0,0) -- (0.8,0);} & \scriptsize implies\\
\tikz{\draw[depuses] (0,0) -- (0.8,0);} & \scriptsize uses\\
\tikz{\draw[depdefines] (0,0) -- (0.8,0);} & \scriptsize defines
\end{tabular}
};
\end{tikzpicture}%
}
% =========================================================
% Dependency Figure Legend
% Standard bottom placement legend for dependency maps
% =========================================================

\newcommand{\DependencyFigureLegend}{%
\begin{tikzpicture}[baseline=(current bounding box.center), >=Latex]
  % Row 1: node legend
  \node[depaxiom, minimum width=0.9cm, minimum height=0.32cm] (ax) at (0,0) {};
  \node[right=0.15cm of ax] (axlab) {\scriptsize Axiom};

  \node[depdefinition, minimum width=0.9cm, minimum height=0.32cm, right=0.9cm of axlab] (df) {};
  \node[right=0.15cm of df] (dflab) {\scriptsize Definition};

  \node[deptheorem, minimum width=0.9cm, minimum height=0.32cm, right=0.9cm of dflab] (th) {};
  \node[right=0.15cm of th] (thlab) {\scriptsize Theorem / Proposition / Corollary};

  % Row 2: edge legend
  \draw[depimplies] ($(ax.south)+(0,-0.9)$) -- ++(0.9,0)
    node[right, text=black] {\scriptsize implies};

  \draw[depuses] ($(df.south)+(0,-0.9)$) -- ++(0.9,0)
    node[right, text=black] {\scriptsize uses};

  \draw[depdefines] ($(th.south)+(0,-0.9)$) -- ++(0.9,0)
    node[right, text=black] {\scriptsize defines};
\end{tikzpicture}%
}
% ---------------------------------------------------------
% Minimal template for a dependency figure file
% Save as figure-<n>.tex and \input it from the notes.
% ---------------------------------------------------------
%
% \begin{center}
% \textbf{<Subchapter or Chapter Title> Dependency Map}
% \end{center}
%
% \begin{tikzpicture}[>=Latex]
%   % Nodes
%   \node[depaxiom]      (a1) at (0,0) {Completeness Axiom};
%   \node[deptheorem]    (t1) at (-3,-1.7) {Nested Interval\\Property};
%   \node[deptheorem]    (t2) at ( 3,-1.7) {Existence of $\sqrt{a}$};
%   \node[deptheorem]    (t3) at (-3,-3.4) {Archimedean Property};
%
%   % Edges
%   \draw[depimplies] (a1) -- node[depedgelab, pos=0.55, left] {implies} (t1);
%   \draw[depimplies] (a1) -- node[depedgelab, pos=0.55, right] {implies} (t2);
%   \draw[depuses]    (t1) -- node[depedgelab, pos=0.50, right] {uses} (t3);
%
%   % Legend (place manually where it fits)
%   \node[anchor=north east] at (6.0,0.2) {\DependencyLegend};
% \end{tikzpicture}

% =========================================================
% Topic-level pedagogical containers
% =========================================================

\newtcolorbox{topicbox}[1]{
  enhanced,
  breakable,
  colback=white,
  colframe=white,
  arc=2pt,
  boxrule=0pt,
  title={\textbf{#1}},
  fonttitle=\bfseries,
  coltitle=black,
  colbacktitle=white,
  boxed title style={boxrule=0pt, colframe=white, colback=white},
  attach boxed title to top left={xshift=0pt, yshift*=-\tcboxedtitleheight/2},
  frame hidden,
  left=0pt, right=0pt, top=8pt, bottom=4pt,
  before skip=8pt, after skip=8pt
}

\NewDocumentEnvironment{exposition}{}%
  {\par\smallskip\noindent\ignorespaces}%
  {\par\smallskip}

\newcommand{\NoLocalDependencies}{}


% ============================================================
% Formal mathematical content boxes
% ============================================================

\newtcolorbox{definitionbox}[1]{
  colback=defbox,
  colframe=defborder,
  arc=2pt,
  left=6pt,
  right=6pt,
  top=4pt,
  bottom=4pt,
  title={\small\textbf{#1}},
  fonttitle=\small\bfseries
}

\newtcolorbox{axiombox}[1]{
  colback=axiombox,
  colframe=axiomborder,
  arc=2pt,
  left=6pt,
  right=6pt,
  top=4pt,
  bottom=4pt,
  title={\small\textbf{#1}},
  fonttitle=\small\bfseries
}

\newtcolorbox{theorembox}[1]{
  colback=thmbox,
  colframe=thmborder,
  arc=2pt,
  left=6pt,
  right=6pt,
  top=4pt,
  bottom=4pt,
  title={\small\textbf{#1}},
  fonttitle=\small\bfseries
}

\newtcolorbox{lemmabox}[1]{
  colback=lembox,
  colframe=lemborder,
  arc=2pt,
  left=6pt,
  right=6pt,
  top=4pt,
  bottom=4pt,
  title={\small\textbf{#1}},
  fonttitle=\small\bfseries
}

\newtcolorbox{propositionbox}[1]{
  colback=propbox,
  colframe=propborder,
  arc=2pt,
  left=6pt,
  right=6pt,
  top=4pt,
  bottom=4pt,
  title={\small\textbf{#1}},
  fonttitle=\small\bfseries
}

\newtcolorbox{corollarybox}[1]{
  colback=corbox,
  colframe=corborder,
  arc=2pt,
  left=6pt,
  right=6pt,
  top=4pt,
  bottom=4pt,
  title={\small\textbf{#1}},
  fonttitle=\small\bfseries
}


% ============================================================
% House colors: navigation / breadcrumbs
% ============================================================
\definecolor{breadcrumb}{HTML}{F2F2F2}
\definecolor{breadcrumbborder}{HTML}{B8B8B8}
\definecolor{breadcrumbtitle}{HTML}{333333}

\providecolor{amber}{HTML}{F59E0B}

% ============================================================
% Breadcrumb navigation + stub status  (deterministic chrome)
% Palette (breadcrumb / breadcrumbborder / breadcrumbtitle) is defined in boxes.tex.
% Emitted by the breadcrumb generator; never hand-rolled, never LLM-drawn.
% Requires: tcolorbox, etoolbox (\ifblank), amsmath (\text); amber color.
% \input this from the preamble after boxes.tex.
% ============================================================
\tcbset{
  breadcrumbstyle/.style={
    colback=breadcrumb, colframe=breadcrumbborder, coltitle=breadcrumbtitle,
    fonttitle=\bfseries\small, arc=2pt, boxrule=0.6pt,
    left=6pt, right=6pt, top=4pt, bottom=4pt,
    before skip=6pt, after skip=8pt,
  }
}

% \breadcrumb{<chapter subject>}{<prior display>}{<current display>}{<next display>}
% Empty prior/next are omitted together with their arrows. ASCII-only content.
\NewDocumentCommand{\breadcrumb}{m m m m}{%
  \begin{tcolorbox}[breadcrumbstyle, title={#1}]%
    \centering$%
      \ifblank{#2}{}{\text{#2}\;\to\;}%
      \textbf{#3}%
      \ifblank{#4}{}{\;\to\;\text{#4}}%
    $%
  \end{tcolorbox}%
}

% \stubstatus  -- appended after the breadcrumb on planned (stub) chapters.
\NewDocumentCommand{\stubstatus}{}{%
  \begin{tcolorbox}[breadcrumbstyle, colback=amber!10, colframe=amber]%
    \textbf{Status:} Planned%
  \end{tcolorbox}%
}

% ============================================================
% Chapter roadmap exposition (replaces the verbose "Structural Roadmap").
% Simple, templated: where the development moves from -> to, and what it studies.
% \chapterroadmap{<prior topic>}{<next topic>}{<topics studied, prose or list>}
% prior/next are the registry neighbours (same source as \breadcrumb); the topic
% list is the chapter's section inventory -- so this can be emitted deterministically.
% ============================================================
\NewDocumentCommand{\chapterroadmap}{m m m}{%
  \begin{exposition}%
  This chapter moves the development from #1 to #2. To get there, it studies: #3.%
  \end{exposition}%
}

% ============================================================
% Colours — front matter
% ============================================================
\definecolor{deepblue}{RGB}{15, 40, 90}
\definecolor{midblue}{RGB}{40, 90, 160}
\definecolor{softgold}{RGB}{180, 140, 40}
\definecolor{lightgray}{RGB}{240, 240, 245}
\definecolor{stageA}{RGB}{20,  80, 140}
\definecolor{stageB}{RGB}{60, 130, 80}
\definecolor{stageC}{RGB}{160, 60, 40}
\definecolor{stageD}{RGB}{120, 50, 140}

% ============================================================
% Colours — mathematical content
% ============================================================
\definecolor{indexblue}  {HTML}{1a3a6b}
\definecolor{natblue}    {HTML}{2e6db4}
\definecolor{realorange} {HTML}{c45c13}
\definecolor{sigmagreen} {HTML}{1a6b3a}
\definecolor{lightblue}  {HTML}{dce8f7}
\definecolor{lightorange}{HTML}{fde8d8}
% Result boxes - one blue family, descending visual weight
\definecolor{thmbox}     {HTML}{f0f6ff}
\definecolor{thmborder}  {HTML}{2f6fae}
\definecolor{propbox}    {HTML}{f4f8ff}
\definecolor{propborder} {HTML}{5f86c7}
\definecolor{lembox}     {HTML}{f7faff}
\definecolor{lemborder}  {HTML}{7fa0d4}
\definecolor{corbox}     {HTML}{fafcff}
\definecolor{corborder}  {HTML}{a8bedf}
% Axiom box - warm amber (unproven primitives)
\definecolor{axiombox}   {HTML}{fff7ed}
\definecolor{axiomborder}{HTML}{c47a2c}
% Definition box - soft teal (first-appearance definitions)
\definecolor{defbox}     {HTML}{f0faf8}
\definecolor{defborder}  {HTML}{3a9a82}

% ============================================================
% Colours — structural presentation roles
% ============================================================
\definecolor{signaturebox}       {HTML}{fff7ed}
\definecolor{signatureborder}    {HTML}{b7791f}
\definecolor{structurebox}       {HTML}{eff6ff}
\definecolor{structureborder}    {HTML}{2f6fae}
\definecolor{structurenewbox}    {HTML}{f0fdf4}
\definecolor{structurenewborder} {HTML}{2f855a}
\definecolor{structuretheorybox} {HTML}{faf5ff}
\definecolor{structuretheoryborder}{HTML}{7b3fb2}
\definecolor{structureusesbox}   {HTML}{f7f4ef}
\definecolor{structureusesborder}{HTML}{8a735a}
\definecolor{classificationbox}  {HTML}{f8fafc}
\definecolor{classificationborder}{HTML}{64748b}
\definecolor{syntacticbox}       {HTML}{f5f3ff}
\definecolor{syntacticborder}    {HTML}{6d5bd0}
\definecolor{formationbox}       {HTML}{f8fafc}
\definecolor{formationborder}    {HTML}{475569}
\definecolor{workedexamplebox}   {HTML}{f0fdf4}
\definecolor{workedexampleborder}{HTML}{2f855a}

% ============================================================
% Colours — algebraic structures diagram (algebraic-dag.tex)
% ============================================================
\definecolor{cSet}    {HTML}{4a4a6a}
\definecolor{cBinop}  {HTML}{5a5a7a}
\definecolor{cMagma}  {HTML}{5a4a7a}
\definecolor{cSemi}   {HTML}{6a5a8a}
\definecolor{cMonoid} {HTML}{7a6a9a}
\definecolor{cGroup}  {HTML}{3a52a4}
\definecolor{cAbelian}{HTML}{2a7a9c}
\definecolor{cRing}   {HTML}{a44a2a}
\definecolor{cField}  {HTML}{a4882a}
\definecolor{cVS}     {HTML}{2a7a4b}
\definecolor{arrISA}  {HTML}{3a5a9a}
\definecolor{arrHASA} {HTML}{8a5a20}
\definecolor{arrUSES} {HTML}{2a7a4a}
\definecolor{nodetext}{HTML}{ffffff}
\definecolor{sigtext} {HTML}{ccccdd}
\definecolor{bodytext}{HTML}{eeeeee}
\definecolor{lbltext} {HTML}{3a3a5a}

% ============================================================
% Structural presentation macros
% Extraction markers: signature, structure, classification, proof-vault-url
% ============================================================

% Structural toolkit: compact navigational summary for section-level vocabulary.
\newtcolorbox{toolkitbox}[1]{
  colback=gray!6,
  colframe=gray!40,
  arc=2pt,
  left=6pt, right=6pt, top=4pt, bottom=4pt,
  title={\small\textbf{#1}},
  fonttitle=\small\bfseries,
  after skip=1em
}

% Signature presentation: model-theoretic languages and symbol tables.
\newtcolorbox{signaturebox}[1]{
  enhanced,
  breakable,
  colback=signaturebox,
  colframe=signatureborder,
  coltitle=black,
  colbacktitle=signatureborder!18!white,
  fonttitle=\bfseries,
  title={#1},
  boxrule=0.8pt,
  arc=2pt,
  left=8pt, right=8pt, top=6pt, bottom=6pt,
  before skip=8pt, after skip=8pt
}

\newenvironment{signaturetable}{%
    \par\smallskip\noindent
    \tabularx{\linewidth}{@{}>{\bfseries}p{0.18\linewidth}p{0.17\linewidth}p{0.13\linewidth}X@{}}
    \toprule
    Symbol & Kind & Arity & Intended Reading\\
    \midrule
  }{%
    \bottomrule
    \endtabularx
    \par\smallskip
  }

% Syntactic objects: variables, terms, formulas, sentences, and formation rules.
\newtcolorbox{syntacticbox}[1]{
  enhanced,
  breakable,
  colback=syntacticbox,
  colframe=syntacticborder,
  coltitle=black,
  colbacktitle=syntacticborder!16!white,
  fonttitle=\bfseries,
  title={#1},
  boxrule=0.8pt,
  arc=2pt,
  left=8pt, right=8pt, top=6pt, bottom=6pt,
  before skip=8pt, after skip=8pt
}

\newtcolorbox{formationrules}[1]{
  enhanced,
  breakable,
  colback=formationbox,
  colframe=formationborder,
  coltitle=black,
  colbacktitle=formationborder!14!white,
  fonttitle=\bfseries\small,
  title={#1},
  boxrule=0.6pt,
  arc=2pt,
  left=7pt, right=7pt, top=5pt, bottom=5pt,
  before skip=6pt, after skip=8pt
}

\newtcolorbox{workedexamples}[1]{
  enhanced,
  breakable,
  colback=workedexamplebox,
  colframe=workedexampleborder,
  coltitle=black,
  colbacktitle=workedexampleborder!14!white,
  fonttitle=\bfseries\small,
  title={#1},
  boxrule=0.6pt,
  arc=2pt,
  left=7pt, right=7pt, top=5pt, bottom=5pt,
  before skip=6pt, after skip=8pt
}

% Structure blueprint: compact inheritance and dependency summaries.
\newtcolorbox{structureblueprint}[1]{
  enhanced,
  breakable,
  colback=structurebox,
  colframe=structureborder,
  coltitle=white,
  colbacktitle=structureborder,
  fonttitle=\bfseries,
  title={#1},
  boxrule=0.8pt,
  arc=2pt,
  left=8pt, right=8pt, top=6pt, bottom=6pt,
  before skip=8pt, after skip=8pt
}

\newtcolorbox{structuresection}[1]{
  enhanced,
  breakable,
  colback=classificationbox,
  colframe=classificationborder,
  coltitle=black,
  colbacktitle=classificationborder!15!white,
  fonttitle=\bfseries\small,
  title={#1},
  boxrule=0.5pt,
  arc=2pt,
  left=6pt, right=6pt, top=4pt, bottom=4pt,
  before skip=4pt, after skip=6pt
}

\newtcolorbox{structurenew}{%
  enhanced,
  breakable,
  colback=structurenewbox,
  colframe=structurenewborder,
  coltitle=black,
  colbacktitle=structurenewborder!16!white,
  fonttitle=\bfseries\small,
  title={NEW AT},
  boxrule=0.5pt,
  arc=2pt,
  left=6pt, right=6pt, top=4pt, bottom=4pt,
  before skip=4pt, after skip=6pt
}

\newtcolorbox{structurehasa}{%
  enhanced,
  breakable,
  colback=structurebox,
  colframe=structureborder,
  coltitle=black,
  colbacktitle=structureborder!14!white,
  fonttitle=\bfseries\small,
  title={HAS-A},
  boxrule=0.5pt,
  arc=2pt,
  left=6pt, right=6pt, top=4pt, bottom=4pt,
  before skip=4pt, after skip=6pt
}

\newtcolorbox{structureusesa}{%
  enhanced,
  breakable,
  colback=structureusesbox,
  colframe=structureusesborder,
  coltitle=black,
  colbacktitle=structureusesborder!16!white,
  fonttitle=\bfseries\small,
  title={USES-A},
  boxrule=0.5pt,
  arc=2pt,
  left=6pt, right=6pt, top=4pt, bottom=4pt,
  before skip=4pt, after skip=6pt
}

\newtcolorbox{structuretheory}{%
  enhanced,
  breakable,
  colback=structuretheorybox,
  colframe=structuretheoryborder,
  coltitle=black,
  colbacktitle=structuretheoryborder!14!white,
  fonttitle=\bfseries\small,
  title={AXIOMS / THEORY},
  boxrule=0.5pt,
  arc=2pt,
  left=6pt, right=6pt, top=4pt, bottom=4pt,
  before skip=4pt, after skip=6pt
}

\newtcolorbox{structureinherits}{%
  enhanced,
  breakable,
  colback=structurebox,
  colframe=structureborder,
  coltitle=black,
  colbacktitle=structureborder!14!white,
  fonttitle=\bfseries\small,
  title={INHERITS},
  boxrule=0.5pt,
  arc=2pt,
  left=6pt, right=6pt, top=4pt, bottom=4pt,
  before skip=4pt, after skip=6pt
}

\newtcolorbox{structureclassification}{%
  enhanced,
  breakable,
  colback=classificationbox,
  colframe=classificationborder,
  coltitle=black,
  colbacktitle=classificationborder!15!white,
  fonttitle=\bfseries\small,
  title={CLASSIFICATION},
  boxrule=0.5pt,
  arc=2pt,
  left=6pt, right=6pt, top=4pt, bottom=4pt,
  before skip=4pt, after skip=6pt
}

% Classification cards: compact metadata for mathematical objects.
\newtcolorbox{classificationbox}[1]{
  enhanced,
  breakable,
  colback=classificationbox,
  colframe=classificationborder,
  coltitle=white,
  colbacktitle=classificationborder,
  fonttitle=\bfseries,
  title={#1},
  boxrule=0.8pt,
  arc=2pt,
  left=8pt, right=8pt, top=6pt, bottom=6pt,
  before skip=8pt, after skip=8pt
}

% ============================================================
% Theorem environments
% ============================================================
\theoremstyle{plain}
\newtheorem{theorem}{Theorem}[section]
\newtheorem{lemma}[theorem]{Lemma}
\newtheorem{proposition}[theorem]{Proposition}
\newtheorem{corollary}[theorem]{Corollary}

\theoremstyle{definition}
\newtheorem{definition}[theorem]{Definition}
\newtheorem{axiom}[theorem]{Axiom}
\newtheorem{assumption}[theorem]{Assumption}

\theoremstyle{remark}
\newtheorem{remark}[theorem]{Remark}
\newtheorem{example}[theorem]{Example}
\newtheorem{principle}[theorem]{Principle}

\newcounter{workedexample}[section]
\renewcommand{\theworkedexample}{\thesection.\arabic{workedexample}}

% Unnumbered variants
\theoremstyle{plain}
\newtheorem*{theorem*}{Theorem}
\newtheorem*{lemma*}{Lemma}
\newtheorem*{proposition*}{Proposition}
\newtheorem*{corollary*}{Corollary}

\newtheorem*{translation}{Translation}
\newtheorem*{notation}{Notation}
\newtheorem*{convention}{Convention}

\theoremstyle{definition}
\newtheorem*{definition*}{Definition}
\newtheorem*{axiom*}{Axiom}
\newtheorem*{assumption*}{Assumption}

% dependencies: typesets "Remark (Dependencies)." as an italic head in the
% digital edition. In print edition, swallow the full body so dependency lists
% remain available to source tooling without occupying page space.
\NewDocumentEnvironment{dependencies}{+b}{%
  \ifdefined\LRAPrintEdition
    \ifLRAPrintEdition
    \else
      \par\addvspace{6pt}%
      \noindent{\itshape Remark} (Dependencies).\par%
      \addvspace{2pt}%
      #1%
      \par\addvspace{6pt}%
    \fi
  \else
    \par\addvspace{6pt}%
    \noindent{\itshape Remark} (Dependencies).\par%
    \addvspace{2pt}%
    #1%
    \par\addvspace{6pt}%
  \fi
}{}

\NewDocumentEnvironment{lra-not-visible}{+b}{}{}

% Worked examples are digital-edition learning artifacts. They are numbered,
% labelable, and source-extractable, but excluded from print builds and do not
% participate in theorem numbering.
\NewDocumentEnvironment{workedexample}{O{} +b}{%
  \LRAExcludeFromPrintEditionBegin
  \refstepcounter{workedexample}%
  \par\addvspace{6pt}%
  \noindent{\itshape Example~\theworkedexample
  \if\relax\detokenize{#1}\relax\else\space(#1)\fi.}\par%
  \addvspace{2pt}%
  #2%
  \par\addvspace{6pt}%
  \LRAExcludeFromPrintEditionEnd
}{}

\theoremstyle{remark}
\newtheorem*{remark*}{Remark}
\newtheorem*{example*}{Example}
\newtheorem*{principle*}{Principle}

\newtheorem{consequence}[theorem]{Consequence}
\newtheorem*{consequence*}{Consequence}

% ============================================================
% Three-column proof table column types (requires array package)
% ============================================================
\newcolumntype{T}{p{0.06\textwidth}}
\newcolumntype{S}{>{\everymath{\displaystyle}$\displaystyle}p{0.44\textwidth}<{$}}
\newcolumntype{J}{p{0.40\textwidth}}

% ============================================================
% Proof-writing macros
% ============================================================
\newcommand{\Given}{\par\noindent\textbf{Given. }\ignorespaces}
\newcommand{\Goal}{\par\noindent\textbf{Goal. }\ignorespaces}
\newcommand{\Strategy}{\par\noindent\textbf{Strategy. }\ignorespaces}
\newcommand{\ProofPlan} {\textbf{Plan.}\ }
\newcommand{\Recall}    {\textbf{Recall.}\ }
\newcommand{\Claim}     {\textbf{Claim.}\ }   % inline bold marker; \begin{claim} = numbered environment
\newcommand{\Wlog}      {\textit{Without loss of generality}}
\newcommand{\Fix}       {\textbf{Fix.}\ }
\newcommand{\Let}       {\textbf{Let.}\ }
\newcommand{\Choose}    {\textbf{Choose.}\ }
\newcommand{\Take}      {\textbf{Take.}\ }
\newcommand{\Note}      {\textbf{Note.}\ }
\newcommand{\ByDef}     {\textbf{By definition,}\ }
\newcommand{\ByAss}     {\textbf{By assumption,}\ }
\newcommand{\ByThm}[1]  {\textbf{By #1,}\ }
\newcommand{\Thus}      {\textbf{Thus,}\ }
\newcommand{\Hence}     {\textbf{Hence,}\ }
\newcommand{\Therefore} {\textbf{Therefore,}\ }
\newcommand{\AsReq}     {\textbf{as required.}}
\newcommand{\checkbox}  {$\square$\ }
\newcommand{\pp}        {\texttt{++}}
\providecommand{\conv}{\operatorname{conv}}
\providecommand{\argmin}{\operatorname*{arg\,min}}
\providecommand{\argmax}{\operatorname*{arg\,max}}

% ============================================================
% Nested numbered claims (reset per proof)
%   Claim 1, Claim 2, ...
%   Subclaim 1.1, 1.2, ...
% ============================================================
\newcounter{claim}
\newcounter{subclaim}[claim]

\renewcommand{\theclaim}{\arabic{claim}}
\renewcommand{\thesubclaim}{\theclaim.\arabic{subclaim}}

\AtBeginEnvironment{proof}{%
  \setcounter{claim}{0}%
  \setcounter{subclaim}{0}%
}

\newenvironment{claim}{%
  \refstepcounter{claim}%
  \par\medskip
  \noindent\textbf{Claim \theclaim.}\ \itshape
}{%
  \par\medskip
}

\newenvironment{subclaim}{%
  \refstepcounter{subclaim}%
  \par\medskip
  \noindent\textbf{Subclaim \thesubclaim.}\ \itshape
}{%
  \par\medskip
}

\newcommand{\ProofStep}[2]{\medskip\noindent\textbf{Step #1\,—\,#2}\quad}
% ============================================================
% Proof annotation tags — FIXED VERSION
% ============================================================
%
% Problem: \marginpar is forbidden inside tabular, tcolorbox, theorem
% environments, and other restricted modes. Calling \prooftag in those
% contexts raises the @marbox / \@marbox undefined error.
%
% Solution: Two separate primitives.
%
%   \prooftag{color}{label}   — margin version (use in proof body prose)
%   \inlinetag{color}{label}  — inline colored badge (use inside tabular T column,
%                               tcolorbox, definition, theorem environments)
%
% All named tag macros (e.g. \tagDU) produce the MARGIN version by default.
% Inside a tabular or box, call the \itag* inline variants instead.
% ============================================================

% ---- Base primitives ----------------------------------------

% Proof tag: inline colored badge.
% Originally used \marginpar but replaced with inline to avoid
% "Floats lost" errors at \clearpage in long documents.
\newcommand{\prooftag}[2]{%
  \textbf{\textcolor{#1}{\scriptsize #2}}%
}

% Inline tag: colored bold badge, zero extra vertical space.
% Safe everywhere: tabular, tcolorbox, theorem bodies, definition boxes.
\newcommand{\inlinetag}[2]{%
  \textbf{\textcolor{#1}{\scriptsize #2}}%
}

% ---- Core structural tags -----------------------------------
%   \tagXX    = margin version  (default, proof prose)
%   \itagXX   = inline version  (tabular T column, boxes)

% DU = Definition Unpacked
\newcommand{\tagDU} {\prooftag{blue!70!black}{DU}}
\newcommand{\itagDU}{\inlinetag{blue!70!black}{DU}}

% TA = Theorem / Axiom Applied
\newcommand{\tagTA} {\prooftag{green!50!black}{TA}}
\newcommand{\itagTA}{\inlinetag{green!50!black}{TA}}

% CS = Case Split
\newcommand{\tagCS} {\prooftag{red!60!black}{CS}}
\newcommand{\itagCS}{\inlinetag{red!60!black}{CS}}

% EX = Existence / witness choice
\newcommand{\tagEX} {\prooftag{magenta!60!black}{EX}}
\newcommand{\itagEX}{\inlinetag{magenta!60!black}{EX}}

% NEG = Negation unpacked (contradiction / contrapositive)
\newcommand{\tagNEG} {\prooftag{black}{NEG}}
\newcommand{\itagNEG}{\inlinetag{black}{NEG}}

% ---- ε–N / tail-logic tags ----------------------------------

% AM = Algebraic Manipulation
\newcommand{\tagAM} {\prooftag{black}{AM}}
\newcommand{\itagAM}{\inlinetag{black}{AM}}

% EA = ε arbitrary (start of universal ε argument)
\newcommand{\tagEA} {\prooftag{black}{EA}}
\newcommand{\itagEA}{\inlinetag{black}{EA}}

% NE = choose N = N(ε)
\newcommand{\tagNE} {\prooftag{black}{NE}}
\newcommand{\itagNE}{\inlinetag{black}{NE}}

% TN = take tail indices ≥ N
\newcommand{\tagTN} {\prooftag{black}{TN}}
\newcommand{\itagTN}{\inlinetag{black}{TN}}

% TR = triangle inequality / norm inequality
\newcommand{\tagTR} {\prooftag{black}{TR}}
\newcommand{\itagTR}{\inlinetag{black}{TR}}

% ---- Completeness / extraction tags -------------------------

% LUB = least-upper-bound / sup step
\newcommand{\tagLUB} {\prooftag{black}{LUB}}
\newcommand{\itagLUB}{\inlinetag{black}{LUB}}

% GLB = greatest-lower-bound / inf step
\newcommand{\tagGLB} {\prooftag{black}{GLB}}
\newcommand{\itagGLB}{\inlinetag{black}{GLB}}

% SS = subsequence selection / extraction
\newcommand{\tagSS} {\prooftag{black}{SS}}
\newcommand{\itagSS}{\inlinetag{black}{SS}}

% ---- Induction tags -----------------------------------------

% IH = inductive hypothesis applied
\newcommand{\tagIH} {\prooftag{purple!70!black}{IH}}
\newcommand{\itagIH}{\inlinetag{purple!70!black}{IH}}

% BC = base case
\newcommand{\tagBC} {\prooftag{teal!70!black}{BC}}
\newcommand{\itagBC}{\inlinetag{teal!70!black}{BC}}

% IS = inductive step
\newcommand{\tagIS} {\prooftag{teal!50!black}{IS}}
\newcommand{\itagIS}{\inlinetag{teal!50!black}{IS}}

% ------------------------------------------------------------
% Removed on purpose (not in ideal minimal set)
% ------------------------------------------------------------
% AM (algebraic manipulation): too granular; prose should carry it.
% UW (universal/witness intro): redundant with TN/EX in practice.
% ------------------------------------------------------------

% ============================================================
% Relation symbols
% ============================================================
\newcommand{\colonequiv}{\vcentcolon\equiv}   % predicate definitional equivalence  P :≡ Q
\newcommand{\coloniff}{\;\vcentcolon\iff\;}   % biconditional form of the above

% ============================================================
% Tolerance / Proximity quantifier macros
%
% Usage in predicate readings:
%
%   \ForallOT{\varepsilon}   --  \forall OutputTolerance(\varepsilon)
%   \ExistsIT{\delta}        --  \exists InputTolerance(\delta)
%   \ForallIT{\delta}        --  \forall InputTolerance(\delta)
%                                (used in uniform continuity: both
%                                 x and u range over inputs)
%
% These make the quantifier-typed reading of epsilon-delta
% definitions compact and semantically explicit.
%
% Full forms (for display when space allows):
%   \forall\,\operatorname{OutputTolerance}(\varepsilon)
%   \exists\,\operatorname{InputTolerance}(\delta)
% ============================================================
\newcommand{\ForallOT}[1]{%
  \forall\,\operatorname{OutputTolerance}(#1)}
\newcommand{\ExistsIT}[1]{%
  \exists\,\operatorname{InputTolerance}(#1)}
\newcommand{\ForallIT}[1]{%
  \forall\,\operatorname{InputTolerance}(#1)}

% ============================================================
% Math operators
% ============================================================
\DeclareMathOperator{\id}{id}
\DeclareMathOperator{\dom}{dom}
\DeclareMathOperator{\cod}{cod}
\DeclareMathOperator{\im}{im}

% ============================================================
% Citation macros
% ============================================================
\newcommand{\srccite}[1]{%
  \textsc{\citeauthor{#1}}%
}
\newcommand{\srcciteauthor}[1]{%
  \textsc{\citeauthor{#1}}%
}
\newcommand{\srccitesection}[2]{%
  \textsc{\citeauthor{#1}}, \textit{#2} [\citeyear{#1}]%
}

% ============================================================
% Proof vault links
% ============================================================
\newcommand{\ProofVaultURL}[1]{%
  \begin{remark*}[Proof vault]
  \href{#1}{Handwritten proof record}
  \end{remark*}%
}

% ============================================================
% Structural metadata markers — no PDF output
% Extraction markers: signature, structure, classification, proof-vault-url
% ============================================================
\newcommand{\LRAProofFor}[1]{}            % proof-file target theorem/proposition/lemma/corollary label
\newcommand{\LRAProofBodyStart}{\mbox{}\par\nobreak\noindent\ignorespaces}
\newcommand{\SignatureMetadata}[1]{}       % future Knowledge Explorer signature node metadata
\newcommand{\StructureMetadata}[1]{}       % future Knowledge Explorer structure blueprint metadata
\newcommand{\ClassificationMetadata}[1]{}  % future Knowledge Explorer classification metadata
\newcommand{\SyntacticMetadata}[1]{}       % future Knowledge Explorer syntactic object metadata

% ============================================================
% Print edition routing
% ============================================================
% Default: complete digital/PDF edition.
% Print edition may be enabled either by:
%   LRA_PRINT_EDITION=1 latexmk -lualatex main.tex
% or by command-line TeX injection:
%   lualatex "\def\LRAPrintEditionRequested{1}\input{main.tex}"
%
% In print edition, theorem-proof vaults, exercise vaults, and capstones are
% not rendered. Source files remain present for local reference and tooling.
\newif\ifLRAPrintEdition
\LRAPrintEditionfalse
\ifdefined\LRAPrintEditionRequested
  \LRAPrintEditiontrue
\fi
\ifdefined\directlua
  \directlua{
    local v = os.getenv("LRA_PRINT_EDITION")
    if v and v ~= "" and v ~= "0" and v:lower() ~= "false" and v:lower() ~= "no" then
      tex.sprint("\\global\\LRAPrintEditiontrue")
    end
  }
\fi

\newcommand{\LRAContentInput}[1]{\input{#1}}
\newcommand{\LRAProofsInput}[1]{\ifLRAPrintEdition\else\input{#1}\fi}
\newcommand{\LRAExercisesInput}[1]{\ifLRAPrintEdition\else\input{#1}\fi}
\newcommand{\LRACapstoneInput}[1]{\ifLRAPrintEdition\else\input{#1}\fi}


% ------------------------------------------------------------
%  Completeness web -- shorthand macros
%  Add to preamble (e.g. macros.tex or main preamble)
% ------------------------------------------------------------

% Plain smallcaps label (matches body text weight)

\newcommand{\AoC}{\hyperref[sec:completeness-web]{\textsc{aoc}}}
\newcommand{\NIP}{\hyperref[sec:nip-proofs]{\textsc{nip}}}
\newcommand{\MCT}{\hyperref[sec:mct-proofs]{\textsc{mct}}}
\newcommand{\BW}{\hyperref[sec:bw-proofs]{\textsc{bw}}}
\newcommand{\CC}{\hyperref[sec:cauchy-proofs]{\textsc{cc}}}

% ============================================================
% Flash macros — all no-ops; produce no PDF output
% Used by scripts/extract-cards.py (and future graph extractors)
% to generate flashcards, theorem-definition explorer nodes,
% logical forms, and dependency metadata.
% ============================================================

% --- Existing core flashcard fields ---
\newcommand{\FlashQ}[1]{}
\newcommand{\FlashA}[1]{}
\newcommand{\FlashHint}[1]{}
\newcommand{\FlashNeg}[1]{}
\newcommand{\FlashImplies}[1]{}
\newcommand{\FlashProof}[1]{}
\newcommand{\FlashUses}[1]{}
\newcommand{\FlashIgnore}{}

% --- New formal / explorer fields ---
\newcommand{\FlashProp}[1]{}        % propositional-style name/form, e.g. ContinuousAt(f,c)
\newcommand{\FlashFormal}[1]{}      % fully formal statement, usually quantified
\newcommand{\FlashNegFormal}[1]{}   % formal negation of the statement/definition
\newcommand{\FlashContra}[1]{}      % contrapositive, when appropriate

% --- New graph / dependency fields ---
\newcommand{\FlashDependsOn}[1]{}   % comma-separated label ids this node depends on
\newcommand{\FlashUsedBy}[1]{}      % comma-separated label ids that consume this node
\newcommand{\FlashPrereq}[1]{}      % softer prerequisite list (can differ from strict proof dependency)
\newcommand{\FlashEquivalentTo}[1]{}% comma-separated equivalent labels / formulations

% --- New classification / search fields ---
\newcommand{\FlashAliases}[1]{}     % alternate names / phrasings
\newcommand{\FlashTags}[1]{}        % comma-separated topical tags
\newcommand{\FlashSymbol}[1]{}      % key symbol or notation introduced by a definition
\newcommand{\FlashCategory}[1]{}    % optional higher grouping: sequences, continuity, compactness, etc.

% --- Optional proof-structure enrichment ---
\newcommand{\FlashProofIdea}[1]{}   % one-line proof idea
\newcommand{\FlashProofType}[1]{}   % contradiction / contrapositive / induction / epsilon-delta / sequential / etc.
\newcommand{\FlashMethod}[1]{}      % named technique(s): squeeze, LUB, BW, MCT, triangle inequality, etc.

% --- Optional note / source enrichment ---
\newcommand{\FlashSourceText}[1]{}  % short bibliographic source text if you want it local to the note
\newcommand{\FlashSection}[1]{}     % local section/subsection string if desired

% ------------------------------------------------------------
% Propositional logic notation
% ------------------------------------------------------------
\providecommand{\False}{\mathsf{False}}
\providecommand{\True}{\mathsf{True}}
\providecommand{\Prop}{\mathsf{Prop}}
\providecommand{\WFF}{\mathsf{WFF}}
\providecommand{\Vars}{\mathsf{Var}}
\providecommand{\NAND}{\mathbin{\uparrow}}
\providecommand{\NOR}{\mathbin{\downarrow}}


% -- Hyperlinks and cross-references -----------------------
\usepackage{hyperref}
\usepackage[nameinlink,capitalize,noabbrev]{cleveref}
\usepackage{bookmark}

% -- cleveref configuration --------------------------------
\crefname{theorem}{Theorem}{Theorems}
\Crefname{theorem}{Theorem}{Theorems}
\crefname{claim}{Claim}{Claims}
\Crefname{claim}{Claim}{Claims}
\crefname{subclaim}{Subclaim}{Subclaims}
\Crefname{subclaim}{Subclaim}{Subclaims}

\hypersetup{
  colorlinks = true,
  linkcolor  = blue!60!black,
  citecolor  = blue!60!black,
  urlcolor   = blue!60!black,
}

% Allow \hyperref inside theorem optional arguments.
\pdfstringdefDisableCommands{%
  \def\hyperref[#1]#2{#2}%
}

% -- Indexes -----------------------------------------------
\usepackage{imakeidx}
\makeindex[name=tech,title={Technique Index}]
\makeindex[name=dep,title={Dependency Index}]

% -- Numbering and TOC depth -------------------------------
\setcounter{secnumdepth}{3}
\setcounter{tocdepth}{3}


\begin{document}
\ExplSyntaxOff
\renewcommand{\familydefault}{\ttdefault}

\tableofcontents
\clearpage

% =========================================================
% Applied and Computational Mathematics
% =========================================================

\part{Applied and Computational Mathematics}
\microtypesetup{expansion=false}
% Book: book-numerical-foundations
% Book router --- no \part; chapters render under the volume's existing \part.
\input{volume-viii/book-numerical-foundations/numerical-analysis/index}
\clearpage
\input{volume-viii/book-numerical-foundations/numerical-pde/index}

\clearpage
% Book: book-applied-methods
% Book router --- no \part; chapters render under the volume's existing \part.
\input{volume-viii/book-applied-methods/computational-calculus/index}
\clearpage
\input{volume-viii/book-applied-methods/geometric-modeling/index}
\clearpage
\input{volume-viii/book-applied-methods/computational-geometry/index}
\clearpage
\input{volume-viii/book-applied-methods/computational-linear-algebra/index}
\clearpage
\input{volume-viii/book-applied-methods/fourier-analysis/index}
\clearpage
\input{volume-viii/book-applied-methods/ordinary-differential-equations/index}
\clearpage
\input{volume-viii/book-applied-methods/partial-differential-equations/index}
\clearpage
\input{volume-viii/book-applied-methods/computational-probability/index}

\microtypesetup{expansion=true}


\clearpage
\bibliographystyle{plainnat}
\bibliography{%
  bibliography/volume-i-foundations,%
  bibliography/volume-ii-number-systems,%
  bibliography/volume-iii-analysis,%
  bibliography/volume-iv-algebra,%
  bibliography/volume-v-topology-geometry,%
  bibliography/volume-vi-computational,%
  bibliography/volume-vii-numerical-approximation,%
  bibliography/volume-viii-logic-foundations,%
  bibliography/general-reference%
}

\printindex[tech]
\printindex[dep]

\end{document}
"
%
% In print edition, theorem-proof vaults, exercise vaults, and capstones are
% not rendered. Source files remain present for local reference and tooling.
\newif\ifLRAPrintEdition
\LRAPrintEditionfalse
\ifdefined\LRAPrintEditionRequested
  \LRAPrintEditiontrue
\fi
\ifdefined\directlua
  \directlua{
    local v = os.getenv("LRA_PRINT_EDITION")
    if v and v ~= "" and v ~= "0" and v:lower() ~= "false" and v:lower() ~= "no" then
      tex.sprint("\\global\\LRAPrintEditiontrue")
    end
  }
\fi

\newcommand{\LRAContentInput}[1]{\input{#1}}
\newcommand{\LRAProofsInput}[1]{\ifLRAPrintEdition\else\input{#1}\fi}
\newcommand{\LRAExercisesInput}[1]{\ifLRAPrintEdition\else\input{#1}\fi}
\newcommand{\LRACapstoneInput}[1]{\ifLRAPrintEdition\else\input{#1}\fi}


% ------------------------------------------------------------
%  Completeness web -- shorthand macros
%  Add to preamble (e.g. macros.tex or main preamble)
% ------------------------------------------------------------

% Plain smallcaps label (matches body text weight)

\newcommand{\AoC}{\hyperref[sec:completeness-web]{\textsc{aoc}}}
\newcommand{\NIP}{\hyperref[sec:nip-proofs]{\textsc{nip}}}
\newcommand{\MCT}{\hyperref[sec:mct-proofs]{\textsc{mct}}}
\newcommand{\BW}{\hyperref[sec:bw-proofs]{\textsc{bw}}}
\newcommand{\CC}{\hyperref[sec:cauchy-proofs]{\textsc{cc}}}

% ============================================================
% Flash macros — all no-ops; produce no PDF output
% Used by scripts/extract-cards.py (and future graph extractors)
% to generate flashcards, theorem-definition explorer nodes,
% logical forms, and dependency metadata.
% ============================================================

% --- Existing core flashcard fields ---
\newcommand{\FlashQ}[1]{}
\newcommand{\FlashA}[1]{}
\newcommand{\FlashHint}[1]{}
\newcommand{\FlashNeg}[1]{}
\newcommand{\FlashImplies}[1]{}
\newcommand{\FlashProof}[1]{}
\newcommand{\FlashUses}[1]{}
\newcommand{\FlashIgnore}{}

% --- New formal / explorer fields ---
\newcommand{\FlashProp}[1]{}        % propositional-style name/form, e.g. ContinuousAt(f,c)
\newcommand{\FlashFormal}[1]{}      % fully formal statement, usually quantified
\newcommand{\FlashNegFormal}[1]{}   % formal negation of the statement/definition
\newcommand{\FlashContra}[1]{}      % contrapositive, when appropriate

% --- New graph / dependency fields ---
\newcommand{\FlashDependsOn}[1]{}   % comma-separated label ids this node depends on
\newcommand{\FlashUsedBy}[1]{}      % comma-separated label ids that consume this node
\newcommand{\FlashPrereq}[1]{}      % softer prerequisite list (can differ from strict proof dependency)
\newcommand{\FlashEquivalentTo}[1]{}% comma-separated equivalent labels / formulations

% --- New classification / search fields ---
\newcommand{\FlashAliases}[1]{}     % alternate names / phrasings
\newcommand{\FlashTags}[1]{}        % comma-separated topical tags
\newcommand{\FlashSymbol}[1]{}      % key symbol or notation introduced by a definition
\newcommand{\FlashCategory}[1]{}    % optional higher grouping: sequences, continuity, compactness, etc.

% --- Optional proof-structure enrichment ---
\newcommand{\FlashProofIdea}[1]{}   % one-line proof idea
\newcommand{\FlashProofType}[1]{}   % contradiction / contrapositive / induction / epsilon-delta / sequential / etc.
\newcommand{\FlashMethod}[1]{}      % named technique(s): squeeze, LUB, BW, MCT, triangle inequality, etc.

% --- Optional note / source enrichment ---
\newcommand{\FlashSourceText}[1]{}  % short bibliographic source text if you want it local to the note
\newcommand{\FlashSection}[1]{}     % local section/subsection string if desired

% ------------------------------------------------------------
% Propositional logic notation
% ------------------------------------------------------------
\providecommand{\False}{\mathsf{False}}
\providecommand{\True}{\mathsf{True}}
\providecommand{\Prop}{\mathsf{Prop}}
\providecommand{\WFF}{\mathsf{WFF}}
\providecommand{\Vars}{\mathsf{Var}}
\providecommand{\NAND}{\mathbin{\uparrow}}
\providecommand{\NOR}{\mathbin{\downarrow}}
